\section*{ВВЕДЕНИЕ}
\addcontentsline{toc}{section}{ВВЕДЕНИЕ}

Современные интернет-технологии активно развиваются и проникают во все сферы жизни. Одним из ярких проявлений этого процесса стало появление разных площадок для торговли вещами, в рамках которой люди получают возможность приобретать вещи как в собственность, так и арендовать их у других пользователей. Онлайн-площадки для аренды вещей позволяют эффективно использовать временно невостребованные предметы, экономя ресурсы и материальные средства. Интенсивность развития данного направления не имеет аналогов — от первых досок объявлений до современных платформ с автоматизированным бронированием, системой рейтингов и встроенными чатами.

При разработке онлайн-площадки все стадии реализации проекта — от проектирования базы данных до развертывания серверной части — находятся в единой технологической цепи. Серверная часть является ядром любой современной веб-платформы, обеспечивая хранение данных, выполнение бизнес-логики и взаимодействие с клиентскими приложениями через программные интерфейсы.

Современные компании, видя, как развиваются информационные технологии, стремятся использовать их для развития своего бизнеса. Онлайн-площадка позволяет не только заявить о себе, но и предоставить пользователям удобный сервис для размещения объявлений, поиска нужных вещей, бронирования и общения. Качественно реализованная серверная часть обеспечивает надежность, безопасность и масштабируемость всего проекта.

\emph{Цель настоящей работы} – разработка серверной части онлайн-площадки для аренды вещей, обеспечивающей хранение данных, бизнес-логику и взаимодействие с клиентской частью. Для достижения поставленной цели необходимо решить \emph{следующие задачи:}
\begin{itemize}
\item провести анализ предметной области и существующих решений в сфере аренды вещей;
\item спроектировать базу данных, выделить основные сущности и установить связи между ними;
\item разработать классы и методы, реализующие бизнес-логику приложения;
\item создать программный интерфейс для взаимодействия с клиентской частью и модулем коммуникации.
\item развернуть и настроить серверную инфраструктуру
\end{itemize}

\emph{Структура и объем работы.} Отчет состоит из введения, 4 разделов основной части, заключения, списка использованных источников, 2 приложений. Текст выпускной квалификационной работы равен \formbytotal{lastpage}{страниц}{е}{ам}{ам}.

\emph{Во введении} сформулирована цель работы, поставлены задачи разработки, описана структура работы, приведено краткое содержание каждого из разделов.

\emph{В первом разделе} проводится анализ предметной области аренды вещей, рассматриваются существующие онлайн-площадки, определяются функциональные требования к разрабатываемой системе.

\emph{Во втором разделе} формулируются требования к серверной части, включая требования к базе данных, программному интерфейсу, производительности и безопасности.

\emph{В третьем разделе} на стадии технического проектирования представлены проектные решения для серверной части.

\emph{В четвертом разделе} описывается реализация серверной части, приводится описание разработанных классов и методов, рассматриваются вопросы развертывания и настройки серверной инфраструктуры, проводится тестирование разработанной системы.

В заключении излагаются основные результаты работы, полученные в ходе разработки.

В приложении А представлен графический материал.
В приложении Б представлены фрагменты исходного кода. 
